\chapter{Outer 同步与有限 inner 的动力学}

\section{Hard sync 的精确动力学}

先考虑每个 history 的 teacher 分布 $T$ 不随 outer 改变，并且 PMD-Part、未线性化 BPO 或 MDPO 的 inner 精确收敛到式~\eqref{eq:token-target} 的 Gibbs 分布。hard sync 定义为
\[
 \mu_k=\pi_k,
 \qquad
 \pi_{k+1}=\frac{\pi_k^{1-\eta}T^\eta}{Z_k}.
\]

\begin{theorem}[Hard-sync geometric interpolation]
\label{thm:hard-sync}
若 $0<\eta\leq1$ 且初始策略与 $T$ 互相绝对连续，则
\begin{equation}\label{eq:hard-closed}
 \pi_k(a\mid s)=
 \frac{\pi_0(a\mid s)^{(1-\eta)^k}
 T(a\mid s)^{1-(1-\eta)^k}}
 {\sum_b\pi_0(b\mid s)^{(1-\eta)^k}T(b\mid s)^{1-(1-\eta)^k}}.
\end{equation}
并且
\[
 \log\frac{\pi_k(a\mid s)}{T(a\mid s)}
 =(1-\eta)^k\log\frac{\pi_0(a\mid s)}{T(a\mid s)}-\log C_k(s),
\]
其中 $C_k(s)$ 是归一化常数。$\eta=1$ 时一个 exact inner 将策略直接变成 teacher；$0<\eta<1$ 时每一轮只消除比例 $\eta$ 的 log-density mismatch。
\end{theorem}

\begin{proof}
归纳证明。$k=0$ 显然成立。假设式~\eqref{eq:hard-closed} 成立，则
\[
 \pi_{k+1}\propto\pi_k^{1-\eta}T^\eta
 \propto \pi_0^{(1-\eta)^{k+1}}
 T^{(1-\eta)(1-(1-\eta)^k)+\eta}.
\]
指数合并为 $1-(1-\eta)^{k+1}$，归一化后得到结论。\hfill$\square$
\end{proof}

\begin{corollary}[Outer contraction]
在任意固定 $s$ 上，unnormalized log-ratio 的 contraction factor 是 $1-\eta$。若 log-ratio 初始振幅不超过 $B$，则归一化后 total variation 满足
\[
 \TV(\pi_k,T)\leq\frac12\left(e^{2B(1-\eta)^k}-1\right).
\]
该界说明小 $\eta$ 不是免费加稳定性：它把每轮变化缩小，同时把达到相同 teacher proximity 所需的 outer steps 增加到 $O(\eta^{-1}\log(1/\varepsilon))$。
\end{corollary}

\section{PMD-Mean 的精确 outer 目标}

PMD-Mean 的精确 inner 不满足前节的 Gibbs 插值递推。固定 $s$，令
\[
G_k(a)=\eta\left(\log\frac{T(a)}{\pi_k(a)}
-\mathbb E_{\pi_k}\left[\log\frac T{\pi_k}\right]\right).
\]
hard sync 的精确递推为
\[
\pi_{k+1}^{\mathrm{Mean}}(a)
=\pi_k(a)\exp\left(G_k(a)-W_0(\lambda_k e^{G_k(a)})\right),
\]
其中 $\lambda_k$ 每轮由归一化重新确定。action-dependent 的 Lambert-$W$ 项阻止了前节的指数合并，因此不能把 $(1-\eta)^k$ 的闭式递推或 $\eta=1$ 的一步 teacher 结论用于 PMD-Mean。

teacher 仍是 PMD-Mean 的固定点：当 $\pi_k=T$ 时，$G_k=0$、$\lambda_k=0$，所以 $\pi_{k+1}=T$。具有 teacher 固定点与每一轮都恢复 Gibbs inner 解是不同结论，后者对 PMD-Mean 不成立。

\section{有限 inner 与 stale target}

本节先以 Gibbs inner 方法为对象。实际第 $k$ 轮只运行 $M_k$ 个 inner steps，得到 $\widehat\pi_{k+1}$。假设在每个 history 上有 log-density 误差
\[
 \left|\log\frac{\widehat\pi_{k+1}(a\mid s)}{\pi_{k+1}^+(a\mid s)}\right|\leq\varepsilon_k
\]
在归一化前成立。则归一化后
\begin{equation}\label{eq:finite-inner-tv}
 \TV(\widehat\pi_{k+1},\pi_{k+1}^+)
 \leq \frac12(e^{2\varepsilon_k}-1).
\end{equation}
几何插值在 log-ratio 振幅上收缩，不能据此前的推论直接宣称 TV 每步以 $1-\eta$ 收缩。定义
\[
b_k=\max_a\log\frac{\pi_k(a\mid s)}{T(a\mid s)}
-\min_a\log\frac{\pi_k(a\mid s)}{T(a\mid s)}.
\]
精确 Gibbs 更新把振幅变成 $(1-\eta)b_k$；上述逐 action log-density 误差最多再增加 $2\varepsilon_k$。因此
\begin{equation}\label{eq:outer-error-recursion}
b_{k+1}\le(1-\eta)b_k+2\varepsilon_k,
\end{equation}
归纳展开得到
\[
b_K\le(1-\eta)^K b_0
+2\sum_{j=0}^{K-1}(1-\eta)^{K-1-j}\varepsilon_j.
\]
归一化分布的 log-ratio 最小值不大于零，最大值不小于零，故 $|\log(\pi_K/T)|\le b_K$。由 $\TV(\pi_K,T)=\frac12\mathbb E_T[|\pi_K/T-1|]$ 得
\[
\TV(\pi_K,T)\le\min\left\{1,\frac12(e^{b_K}-1)\right\}.
\]
若 $\varepsilon_j\le\bar\varepsilon$，振幅的稳态上界为 $2\bar\varepsilon/\eta$。这给出有限 inner 误差的几何累积关系；对 PMD-Mean 需要另行分析其 Lambert-$W$ 更新，不能套用该递推。

\section{Soft sync 的耦合动力学}

若 anchor 使用 EMA，令
\[
 \mu_{k+1}=\lambda\mu_k+(1-\lambda)\widehat\pi_{k+1}
\]
表示概率空间中的近似 soft sync。对 PMD-Part、BPO 和 MDPO，精确 inner 目标是
\[
 \widehat\pi_{k+1}\propto\mu_k^{1-\eta}T^\eta.
\]
忽略归一化扰动并在线性 log-density 坐标中写 $z_k=\log\mu_k-\log T$，得到近似递推
\[
 z_{k+1}\approx[\lambda+(1-\lambda)(1-\eta)]z_k
 =[1-(1-\lambda)\eta]z_k.
\]
对 PMD-Mean，必须把上述 Gibbs target 替换为 Lambert-$W$ target，不能沿用这个线性化系数。对 Gibbs inner 方法，EMA 同时把 outer contraction 从 $1-\eta$ 改成 $1-(1-\lambda)\eta$。当 $\lambda$ 接近 $1$，anchor 漂移非常慢，训练稳定但 teacher tracking 速度下降；这不是 inner 学习率的效果，而是 outer target 的真实动力学。

\section{访问测度固定的含义}

本章所有递推都是逐 history 的 token distribution 递推。真实 LLM 中 history 集合由 $\rho_k$ 产生；只要 inner 内固定该集合，以上结论适用于 conditional token target。outer rollout 更新后，$\rho_{k+1}$ 改变，必须额外记录 coverage drift：
\[
 \Delta_k=\TV(\rho_{k+1},\rho_k).
\]
若 token loss 在所有 history 上有上界 $C$，则访问测度改变给目标梯度带来至多 $2C\Delta_k$ 的偏差。因而“固定测度下求解 token 子问题”和“outer 重新 rollout”是两个严格分开的步骤。
